% Generated draft: requires mathematical and source review.
\documentclass[11pt]{article}
\usepackage[margin=25mm]{geometry}
\usepackage{fontspec}
\setmainfont{Latin Modern Roman}
\usepackage{amsmath,amssymb,mathtools,unicode-math}
\setmathfont{Latin Modern Math}
\usepackage{enumitem,needspace,xurl,hyperref}
\setlength{\parindent}{0pt}
\setlength{\parskip}{5pt}
\setlength{\emergencystretch}{3em}
\newcommand{\N}{\mathbb{N}}
\newcommand{\Z}{\mathbb{Z}}
\newcommand{\Q}{\mathbb{Q}}
\newcommand{\R}{\mathbb{R}}
\newcommand{\C}{\mathbb{C}}
\newcommand{\F}{\mathbb{F}}
\newcommand{\A}{\mathbb{A}}
\newcommand{\PP}{\mathbb{P}}
\newcommand{\E}{\mathbb{E}}
\newcommand{\eps}{\varepsilon}
\newcommand{\dd}{\,\mathrm{d}}
\begin{document}
\section*{UnsolvedMath ID: 700556}
\textbf{Source problem code:} catalogue entry
\subsection{Definitions and mathematical statement}
Let \(p\) be a prime, let \(G\) be a finite group, and let \(B\) be an ordinary \(p\)-block of \(G\), with defect group \(D=D(B)\). Here an ordinary \(p\)-block means a block in the usual modular representation theoretic partition of the irreducible complex characters of \(G\). The source leaves the choice of prime and this interpretation of ``block'' implicit; the statements below use the standard interpretation that they range over every prime \(p\), every finite group \(G\), and every ordinary \(p\)-block \(B\) of \(G\).

Write \(\operatorname{Irr}(B)\) for the set of irreducible complex characters belonging to \(B\), and put \(k(B)=|\operatorname{Irr}(B)|\). For a positive integer \(n\), let \(n_p\) denote the largest power of \(p\) dividing \(n\). Write \(|G|_p=p^a\) and \(|D|=p^d\). The height \(h_B(\chi)\) of \(\chi\in\operatorname{Irr}(B)\) is the nonnegative integer determined by
\[
\chi(1)_p=p^{a-d+h_B(\chi)}.
\]
Define
\[
k_0(B)=\bigl|\{\chi\in\operatorname{Irr}(B):h_B(\chi)=0\}\bigr|,
\qquad
D'=\langle x^{-1}y^{-1}xy:x,y\in D\rangle.
\]
A finite group is \(p\)-soluble if it has a finite normal series whose factors are either \(p\)-groups or groups of order coprime to \(p\).

The source presents the following three assertions, universally quantified as above:
\begin{enumerate}
\item[(a)] \textup{(R. Brauer)}
\[
k(B)\leq |D(B)|.
\]
The source reports that this inequality has been proved for \(p\)-soluble groups, citing D. Gluck, K. Magaard, U. Riese, and P. Schmid (2004).
\item[(b)] \textup{(J. B. Olsson)}
\[
k_0(B)\leq |D(B):D(B)'|.
\]
\item[(c)] \textup{(R. Brauer)}
\[
D(B)\text{ is abelian}\quad\Longleftrightarrow\quad k_0(B)=k(B).
\]
Equivalently, the defect group is abelian if and only if every irreducible complex character in the block has height zero. The supplied discussion labels this assertion solved and attributes its proof to G. Malle, G. Navarro, A. A. Schaeffer Fry, and P. H. Tiep (2024).
\end{enumerate}
The supplied registry labels the combined problem \texttt{partially\_solved}. This label and the reported results are supplied source assessments, not independent mathematical or literature verification.
\subsection{Short English statement}
For every prime \(p\) and every block of irreducible complex characters of a finite group, the problem asks for two bounds on character counts in terms of the block's defect group, and for the equivalence between an abelian defect group and all characters in the block having height zero. The source reports (a) proved for \(p\)-soluble groups and (c) proved in full; it labels the combined record partially solved.
\subsection{Sources}
\begin{itemize}
\item Local UnsolvedMath registry, supplied complete record with permanent ID \(700556\), upstream ID \(700556\), exact source problem code \texttt{catalogue entry}, and title \emph{Kourovka Notebook Problem 9.23}. Supplied upstream revision: \texttt{d371ee917ac9647fbb06cd6fc529deec6f3dda82}.
\item \emph{The Kourovka Notebook (Unsolved Problems in Group Theory)}, original problem 9.23, 9th issue (1984); supplied edition: 21st issue (2026), arXiv:1401.0300v48. The supplied collection identifies E. I. Khukhro and V. D. Mazurov as editors. Proposed by V. D. Mazurov in 1984. Assertions (a) and (c) are attributed to R. Brauer, and assertion (b) to J. B. Olsson.
\item D. Gluck, K. Magaard, U. Riese, and P. Schmid, \emph{J. Algebra} \textbf{279} (2004), 694--719. Supplied reference for assertion (a) for \(p\)-soluble groups; the record does not supply the article title.
\item G. Malle, G. Navarro, A. A. Schaeffer Fry, and P. H. Tiep, \emph{Ann. Math.} \textbf{200} (2024), 557--608. Supplied reference for the reported proof of assertion (c); the record does not supply the article title.
\item Supplied status review dated \texttt{2026-10-16}, with evidence classified as \texttt{maintained\_source}, assigns \texttt{partially\_solved} and reports that assertion (c) is marked solved in the 21st issue, described in that review as October 2026, arXiv:1401.0300v48. This is supplied dated curation, not an independently verified literature assessment.
\end{itemize}
\end{document}
